\section{Procedure}
\begin{enumerate}
  \item The diameter of the spherical brass bob was measured in three orthogonal directions using the digital vernier caliper to calculate its average radius, $r$.
  \item The nylon string was secured to the suspension clamp at the top and attached to the bob at the bottom.
  \item The distance from the bottom of the clamp to the top of the bob was measured with the meter stick to find the string length, $L_{\text{string}}$. The total pendulum length was calculated as $L = L_{\text{string}} + r$.
  \item The photogate was aligned at the lowest point of the pendulum swing path (equilibrium position) and configured in pendulum mode, which measures the time elapsed between three consecutive gate blocks (one complete cycle).
  \item The pendulum bob was displaced from equilibrium by a small angle ($\theta < \SI{10}{\degree}$) and released from rest to ensure stable, planar oscillations.
  \item The period $T$ of a single full oscillation was recorded. This was repeated three times for each pendulum length to compute a mean period.
  \item Steps 3–6 were repeated for five different lengths, varying the total pendulum length from $\SI{0.200}{\meter}$ to $\SI{1.000}{\meter}$ in steps of $\SI{0.200}{\meter}$.
\end{enumerate}

% --- Section: Data & Measurements ---
